In previous sections we presented our algorithm for solving the LWE problem over $\Z_q$ and its variants. In this section we discuss how to solve the problem by decomposing it into multiple polynomial system solving instances over smaller rings (or even finite fields). To provide a direct comparison with the blind application of our algorithm over $\Z_q$, we provide the previous complexity estimates in terms of bit operations. 

Let $B(\log q)$ be a function which bounds the number of bit operations required to perform a single scalar multiplication over $\mathbb{Z}_q$. Using
standard integer arithmetic, $B(\log q) << \log q$, while fast integer arithmetic allows $B(\log q) << M(\log q) \log q \log \log q$, where $M(t)$ is the complexity of multiplying two $t$-bit integers~\cite{DBLP:conf/issac/Storjohann96}.

Let \(l_{\texttt{prob}}\) be the number of variables in the polynomial system corresponding to the problem \texttt{prob} $\in \{\text{LWE}, \rlwe, \mlwe\}$, namely \(l_{\text{LWE}} = n\), \(l_{\rlwe} = N\), and \(l_{\mlwe} = Nn\). The complexity of our algorithm for solving the three LWE variants with bounded errors (e.g. error entries/coefficients drawn from a set of size $d$) is \[\mathscr{C}^{(q)}_{alg} = \bigO\left(l_{\texttt{prob}}^{\max(3(d-1), d\omega)}B(\log q)\right)\] bit operations. 


\textbf{Mapping to smaller rings.} Consider the ring \(\Z_q\) where \(q\) is a composite integer. 
We can factor \(q\) into its prime power factors as \(q = \prod_{i=1}^r p_i^{e_i}\). By the Chinese Remainder Theorem (CRT), we have an isomorphism:
\[\Z_{q} \cong \prod_{i=1}^r \Z_{p_i^{e_i}}.\]
This means that we can represent elements of \(\Z_{q}\) as tuples of elements in the rings \(\Z_{p_i^{e_i}}\). In particular, we can represent the LWE samples as tuples of samples in the rings \(\Z_{p_i^{e_i}}\). This allows us to apply our algorithm to each component of the tuple separately, effectively solving the LWE problem over the finite fields \(\F_{p_i}\) for each prime factor \(p_i\) of \(q\). By exploiting Hensel's lifting, we can lift the solutions obtained in the finite fields to obtain a solution in the corresponding extension fields \(\F_{p_i^{e_i}}\) and then combine the solutions using the CRT to obtain a solution in \(\Z_{q}\).

\noindent
\textbf{Practical considerations.} While in theory this approach permits to avoid scalar multiplication in the $S$-polynomial computation (the diagonalization would generate monic polynomials), in practice it may not be more efficient than working directly over \(\Z_q\), due to the overhead of the mappings and the fact that the algorithm's complexity would depend on the number of elements in the prime factorization of $q$. Therefore, it is important to carefully analyze the trade-offs and consider the specific parameters of the problem when deciding whether to use this approach. As a result, this consideration could help in identifying possibly weak instances.

When $q = \prod_{i=1}^r p_i$, the algorithm is applied over each \(\F_{p_i}\) and can be slightly more efficient (depending on the number of primes in the factorization). In particular, the running time of the algorithm (\emph{not the asymptotic complexity}) is improved by replacing the Smith Normal Form with the Reduced Row Echelon Form which leads to monic polynomials and thus avoids the need of scalar multiplications in the $S$-polynomial computation and reduction. 

When, instead, $q = \prod_{i=1}^r p_i^{e_i}$ with $e_i > 1$ for some $i$, applying the algorithm over \(\F_{p_i}\) and then lifting to higher dimensions may introduce additional overhead and complexity which could become a bottleneck in the overall algorithm. In this case, a better approach is to directly work over the rings \(\Z_{p_i^{e_i}}\). If \(\sum_{i=1}^r \log_2(p_i^{e_i}) << \log_2(q)\), due to the smaller size of the rings \(\Z_{p_i^{e_i}}\) compared to \(\Z_q\), repeating the algorithm $r$ times could be more efficient than working directly over \(\Z_q\).
More generically, considering the complexity of the final CRT, given by \(\bigO(\log_2^2(q))\)~\cite{DBLP:books/daglib/0018101}, an advantage of this approach arises when 
\begin{equation*}\sum_{i=1}^r \mathscr{C}^{(p_i^{e_i})}_{alg} + \mathscr{C}_{CRT} < \mathscr{C}^{(q)}_{alg} \mapsto
    \sum_{i=1}^r B(\log p_i^{e_i})  < B(\log q).
\end{equation*}